\section{Solution Design}
\label{sec:solution}

This section presents the architecture of OpenPA, mapping each design decision to the sub-problems defined in Section~\ref{sec:problem}.

\subsection{System Architecture Overview}

OpenPA follows a three-tier architecture:

\begin{itemize}[leftmargin=*]
  \item \textbf{Frontend}: A Next.js 16 \cite{nextjs2024} web application providing the chat interface, settings management, and real-time streaming via Server-Sent Events (SSE).
  \item \textbf{Backend}: A FastAPI \cite{fastapi2024} server hosting the agent loop, LLM provider abstraction, MCP tool registry, RAG system, and REST API.
  \item \textbf{Data Layer}: SQLite for relational data (users, conversations, credentials) and ChromaDB \cite{chromadb2024} for vector embeddings (semantic memory and conversation search).
\end{itemize}

Figure~\ref{fig:architecture} illustrates the high-level architecture and data flow between components.

\begin{figure}[H]
\centering
\resizebox{\textwidth}{!}{%
\begin{tikzpicture}[
    % Base styles
    node distance=0.6cm and 0.8cm,
    >=stealth,
    % Box styles
    base/.style={
        draw, rounded corners=4pt, minimum height=0.9cm,
        font=\sffamily\small, align=center, thick
    },
    % Layer styles
    user/.style={base, fill=gray!12, draw=gray!60, minimum width=2.8cm},
    frontend/.style={base, fill=blue!8, draw=blue!50, minimum width=2.8cm},
    backend/.style={base, fill=green!8, draw=green!50!black!40, minimum width=2.8cm},
    toolns/.style={base, fill=green!5, draw=green!40!black!30, minimum width=2.4cm,
                   minimum height=0.75cm, font=\sffamily\scriptsize},
    datalayer/.style={base, fill=violet!10, draw=violet!50, minimum width=2.8cm},
    external/.style={base, fill=orange!10, draw=orange!50!black!40, minimum width=2.8cm},
    % Group box style
    groupbox/.style={draw=#1!40, rounded corners=6pt, thick,
                     fill=#1!3, inner sep=8pt},
    % Arrow styles
    dataflow/.style={->, thick, color=gray!70},
    biflow/.style={<->, thick, color=gray!70},
    sseflow/.style={->, thick, color=blue!50, densely dashed},
]

% ============================================================
% USER
% ============================================================
\node[user, minimum width=3.2cm, minimum height=1cm]
    (user) {\textbf{User}\\{\scriptsize Browser / Client}};

% ============================================================
% FRONTEND LAYER
% ============================================================
\node[frontend, below=1.0cm of user, xshift=-2.2cm] (chatui) {Chat UI\\{\scriptsize SSE Streaming}};
\node[frontend, right=0.5cm of chatui] (settings) {Settings\\{\scriptsize API Key Mgmt}};
\node[frontend, right=0.5cm of settings] (oauth-fe) {OAuth\\{\scriptsize Connections}};

% Frontend group box
\begin{scope}[on background layer]
    \node[groupbox=blue, fit=(chatui)(settings)(oauth-fe)] (fe-group) {};
\end{scope}
\node[font=\sffamily\bfseries\small, blue!60!black, fill=white, inner sep=2pt]
    at (fe-group.north) {Frontend (Next.js 16)};

% ============================================================
% BACKEND LAYER
% ============================================================
% Row 1: Core components
\node[backend, below=1.4cm of chatui, xshift=-0.3cm] (agent) {Agent Loop\\{\scriptsize Plan-then-Execute /}\\{\scriptsize Adaptive}};
\node[backend, right=0.5cm of agent] (llmprov) {LLM Provider\\{\scriptsize Abstraction}\\{\scriptsize (6 providers)}};
\node[backend, right=0.5cm of llmprov] (mcpreg) {MCP Tool\\{\scriptsize Registry}\\{\scriptsize (88 tools)}};
\node[backend, right=0.5cm of mcpreg] (rag) {RAG Engine\\{\scriptsize nomic-embed-text}};

% Row 2: Supporting components
\node[backend, below=0.7cm of agent, xshift=1.5cm] (restapi) {REST API\\{\scriptsize + SSE}};
\node[backend, right=0.5cm of restapi] (oauthmgr) {OAuth\\{\scriptsize Manager}};

% Backend group box
\begin{scope}[on background layer]
    \node[groupbox=green, fit=(agent)(llmprov)(mcpreg)(rag)(restapi)(oauthmgr),
          inner sep=10pt] (be-group) {};
\end{scope}
\node[font=\sffamily\bfseries\small, green!40!black, fill=white, inner sep=2pt]
    at (be-group.north) {Backend (FastAPI)};

% ============================================================
% MCP TOOL NAMESPACES
% ============================================================
\node[toolns, below=1.6cm of restapi, xshift=-2.0cm]
    (comm) {\textbf{Communication}\\Gmail\,\textperiodcentered\,Discord\,\textperiodcentered\,Telegram\\{\tiny (12 tools)}};
\node[toolns, right=0.35cm of comm]
    (dev) {\textbf{Development}\\GitHub\,\textperiodcentered\,Workspace\,\textperiodcentered\,Sandbox\\{\tiny (35 tools)}};
\node[toolns, right=0.35cm of dev]
    (media) {\textbf{Media \& Social}\\Spotify\,\textperiodcentered\,Mastodon\,\textperiodcentered\,YouTube\\{\tiny (16 tools)}};
\node[toolns, right=0.35cm of media]
    (prod) {\textbf{Productivity}\\Calendar\,\textperiodcentered\,RSS\,\textperiodcentered\,Scheduler\,\textperiodcentered\,Weather\\{\tiny (12 tools)}};
\node[toolns, right=0.35cm of prod]
    (intel) {\textbf{Intelligence}\\Memory\,\textperiodcentered\,Web Search\,\textperiodcentered\,Scrape\\{\tiny (12 tools)}};

% MCP group box
\begin{scope}[on background layer]
    \node[groupbox=green!70!black, fit=(comm)(dev)(media)(prod)(intel),
          inner sep=8pt] (mcp-group) {};
\end{scope}
\node[font=\sffamily\bfseries\small, green!30!black, fill=white, inner sep=2pt]
    at (mcp-group.north) {MCP Tool Namespaces (16 namespaces)};

% ============================================================
% DATA LAYER
% ============================================================
\node[datalayer, below=1.5cm of dev, xshift=-0.5cm]
    (sqlite) {\textbf{SQLite}\\{\scriptsize Users, Conversations,}\\{\scriptsize Credentials}};
\node[datalayer, right=1.5cm of sqlite]
    (chroma) {\textbf{ChromaDB}\\{\scriptsize Vector Embeddings}\\{\scriptsize (Memory + Search)}};

% Data layer group box
\begin{scope}[on background layer]
    \node[groupbox=violet, fit=(sqlite)(chroma),
          inner sep=10pt] (data-group) {};
\end{scope}
\node[font=\sffamily\bfseries\small, violet!60!black, fill=white, inner sep=2pt]
    at (data-group.south) {Data Layer};

% ============================================================
% EXTERNAL SERVICES (right side)
% ============================================================
\node[external, right=1.8cm of rag, yshift=0.3cm]
    (llmapis) {\textbf{LLM APIs}\\{\scriptsize Anthropic\,\textperiodcentered\,Google}\\{\scriptsize OpenAI\,\textperiodcentered\,Groq\,\textperiodcentered\,OpenRouter}};
\node[external, below=0.5cm of llmapis]
    (ollama) {\textbf{Ollama}\\{\scriptsize Local Models}};
\node[external, below=0.5cm of ollama]
    (svcapis) {\textbf{Service APIs}\\{\scriptsize Gmail\,\textperiodcentered\,GitHub\,\textperiodcentered\,Spotify}\\{\scriptsize Discord\,\textperiodcentered\,Telegram\,\textperiodcentered\,etc.}};

% External group box
\begin{scope}[on background layer]
    \node[groupbox=orange, fit=(llmapis)(ollama)(svcapis),
          inner sep=10pt] (ext-group) {};
\end{scope}
\node[font=\sffamily\bfseries\small, orange!60!black, fill=white, inner sep=2pt]
    at (ext-group.north) {External Services};

% ============================================================
% ARROWS: User <-> Frontend
% ============================================================
\draw[biflow] (user.south) -- ++(0,-0.35) -| (chatui.north);
\draw[biflow] (user.south) -- ++(0,-0.35) -| (settings.north);
\draw[biflow] (user.south) -- ++(0,-0.35) -| (oauth-fe.north);

% ============================================================
% ARROWS: Frontend <-> Backend
% ============================================================
\draw[sseflow] (chatui.south) -- ++(0,-0.5) -| (restapi.north);
\draw[dataflow] (settings.south) -- ++(0,-0.5) -| (restapi.north west);
\draw[dataflow] (oauth-fe.south) -- ++(0,-0.5) -| (oauthmgr.north);

% ============================================================
% ARROWS: Internal backend
% ============================================================
\draw[dataflow] (restapi.north) -- ++(0, 0.25) -| (agent.south);
\draw[dataflow] (agent.east) -- (llmprov.west);
\draw[dataflow] (agent.south east) -- ++(0.3,-0.1) -| (mcpreg.south);
\draw[dataflow] (agent) -- ++(0,0) |- ([yshift=-0.15cm]rag.west);
\draw[dataflow] (mcpreg.south) -- ++(0,-0.65) -- (mcp-group.north);

% ============================================================
% ARROWS: Backend -> External
% ============================================================
\draw[dataflow] (llmprov.east) -- (llmapis.west);
\draw[dataflow] (rag.east) -- (ollama.west);
\draw[dataflow] (oauthmgr.east) -| ([xshift=-0.3cm]svcapis.west) -- (svcapis.west);

% MCP -> Service APIs
\draw[dataflow] (mcp-group.east) -| ([xshift=-0.8cm]svcapis.south west) -- (svcapis.south west);

% ============================================================
% ARROWS: Backend -> Data Layer
% ============================================================
\draw[dataflow] (restapi.south) -- ++(0,-0.4) -| (sqlite.north);
\draw[dataflow] (rag.south) -- ++(0,-1.8) -| (chroma.north);

% ============================================================
% LEGEND
% ============================================================
\node[anchor=north west, font=\sffamily\scriptsize] at ([xshift=0.3cm, yshift=-0.2cm]data-group.south west) {
    \tikz[baseline=-0.5ex]{\draw[->, thick, gray!70] (0,0) -- (0.6,0);} Data flow \quad
    \tikz[baseline=-0.5ex]{\draw[->, thick, blue!50, densely dashed] (0,0) -- (0.6,0);} SSE stream \quad
    \tikz[baseline=-0.5ex]{\draw[<->, thick, gray!70] (0,0) -- (0.6,0);} Bidirectional
};

\end{tikzpicture}%
}% end resizebox
\caption{High-level architecture of the OpenPA system. The frontend communicates with the FastAPI backend via REST and SSE. The agent loop orchestrates LLM inference, tool execution across 16 MCP namespaces (88 tools), and RAG-based retrieval. Data persists in SQLite and ChromaDB, while external LLM and service APIs are accessed through provider abstractions and OAuth-managed credentials.}
\label{fig:architecture}
\end{figure}


\subsection{MCP-Based Tool Architecture (P1)}

To address the service integration problem, OpenPA adopts Anthropic's Model Context Protocol (MCP) \cite{anthropic2024mcp} as the universal tool interface. Each service is implemented as a self-contained FastMCP \cite{fastmcp2024} server with typed tool definitions.

\begin{lstlisting}[language=Python, caption={Example tool definition (simplified)}]
@mcp.tool()
async def send_message(_user_id: int, to: str, message: str):
    """Send a Telegram message. Auto-resolves contact names."""
    creds = await get_creds(_user_id, "telegram")
    # ... send via Telegram API
\end{lstlisting}

All 16 service servers are mounted into a composite registry with namespace prefixes:

\begin{lstlisting}[language=Python, caption={Tool registry mounting}]
mcp = FastMCP("openpa")
mcp.mount(gmail_mcp, namespace="gmail")      # gmail_send_email, ...
mcp.mount(github_mcp, namespace="github")    # github_create_pr, ...
mcp.mount(workspace_mcp, namespace="ws")     # ws_workspace_run, ...
# ... 12 more services
\end{lstlisting}

This architecture yields 88 tools across 16 namespaces. Each tool receives a \texttt{\_user\_id} parameter injected by the agent loop, ensuring per-user credential isolation. OAuth tokens are stored per-user in the database and fetched on demand via a \texttt{get\_creds(user\_id, service)} helper.

Table~\ref{tab:services} summarises the service integrations.

\begin{table}[H]
\centering
\caption{OpenPA service integrations}
\label{tab:services}
\small
\begin{tabularx}{\textwidth}{l|c|X}
\toprule
\textbf{Service} & \textbf{Tools} & \textbf{Key Capabilities} \\
\midrule
GitHub      & 13 & Repos, PRs, issues, branches, file read/write, notifications \\
Workspace   & 15 & Clone, edit, grep, inspect, run tests/builds, commit, push \\
Memory      &  8 & Preferences, semantic memories, conversation search \\
Gmail       &  4 & Unread, read, send, reply (with search syntax) \\
Mastodon    &  9 & Timelines, trending, search, post, notifications \\
Spotify     &  5 & Play, pause, search, playlists, current track \\
Discord     &  4 & Servers, channels, send/read messages \\
Telegram    &  4 & Contacts, send/read messages, chat listing \\
Calendar    &  3 & List, create, delete events \\
Sandbox     &  7 & Run Python/JS, multi-file test, verify syntax, export files \\
Scheduler   &  3 & Schedule future tool calls, list, cancel \\
Web Scrape  &  3 & Fetch page text, extract tables, extract links \\
RSS         &  5 & Fetch, list, add, remove feeds \\
YouTube     &  2 & Download video, get video info \\
Web Search  &  1 & DuckDuckGo search \\
Weather     &  2 & Current conditions, forecast \\
\midrule
\textbf{Total} & \textbf{88} & \\
\bottomrule
\end{tabularx}
\end{table}

\subsection{Adaptive Agent Loop (P2)}

The agent loop is the core reasoning engine. OpenPA implements a dual-mode execution model:

\subsubsection{Mode A: Plan-Then-Execute}

For multi-step tasks that do \textit{not} involve coding (e.g., ``search the web for coffee roasters, compose a summary, and send it via Telegram''), the agent first generates an execution plan as a JSON array:

\begin{lstlisting}[language=Python, caption={Example execution plan}]
[
  {"step": 1, "tool": "web_search", "args": {"query": "..."}},
  {"step": 2, "tool": "llm_generate", "args": {"prompt": "Summarise..."}, "depends_on": 1},
  {"step": 3, "tool": "telegram_send_message", "args": {"to": "...", "content": ""}, "depends_on": 2}
]
\end{lstlisting}

Steps linked by \texttt{depends\_on} receive the output of the referenced step. A regex-based template replacement system handles various placeholder formats (e.g., \texttt{\{\{step\_2.output\}\}}) that LLMs tend to generate, with a safety net that replaces any remaining \texttt{\{\{...\}\}} templates with the most recent step's result.

\subsubsection{Mode B: Adaptive Loop}

For coding tasks and simple queries, the agent uses an adaptive loop where the LLM decides the next action based on tool results. This is critical for vibe-coding, where test failures require reading the error, fixing the code, and re-testing---a flow that cannot be pre-planned.

\begin{lstlisting}[language=Python, caption={Simplified adaptive loop}]
while iterations < max_iterations:
    response = await provider.chat(messages, tools, system_prompt)
    if response.has_tool_calls:
        for tool_call in response.tool_calls:
            result = await call_tool(tool_call.name, tool_call.arguments)
            # Result appended to conversation; LLM sees it next iteration
    else:
        break  # LLM chose to respond with text; task complete
\end{lstlisting}

The planner explicitly routes coding tasks to the adaptive loop:

\begin{quote}
\textit{``If the request involves CODING, WORKSPACE, or VIBE-CODING, ALWAYS output \texttt{\{"simple": true\}}. These tasks MUST use the adaptive loop so you can react to test failures.''}
\end{quote}

Workspace tasks receive an extended iteration limit (50 vs.\ 30) to accommodate multiple test-fix cycles. Progress reminders are injected every three iterations to prevent model drift.

\subsubsection{Tool Filtering (P5)}

To address the model compatibility problem, OpenPA implements context-aware tool filtering. Tools are categorised into nine semantic groups (core, workspace, github, communication, media, web, social, productivity, sandbox). Keywords in the user's message determine which categories are exposed:

\begin{lstlisting}[language=Python, caption={Tool filtering example}]
# "Add a health check endpoint" -> workspace + github + core = 34 tools
# "Play some music"             -> media + core = 15 tools
# "Who am I?"                   -> core only = 8 tools
\end{lstlisting}

This reduces the tool set from 88 to as few as 8 tools, enabling smaller models (7--9B parameters) to function effectively. When workspace tools are called during execution, the filter dynamically expands to include GitHub and web tools for related operations.

\subsubsection{Error Recovery}

Inspired by the Pi coding agent \cite{badlogic2025pi}, tool errors are returned as structured JSON messages rather than exceptions:

\begin{lstlisting}[language=Python, caption={Error-as-tool-result pattern}]
{"error": true, "tool": "github_create_pr",
 "message": "Validation Failed: base branch 'main' not found",
 "hint": "Read the error message, fix the issue, and try again."}
\end{lstlisting}

This allows the LLM to read the error, reason about the cause, and take corrective action---such as changing the base branch from \texttt{main} to \texttt{master}.

\subsection{RAG-Based Personalisation (P3)}

OpenPA's personalisation system uses Retrieval-Augmented Generation to build a persistent understanding of each user.

\subsubsection{Architecture}

\begin{itemize}[leftmargin=*]
  \item \textbf{Embedding model}: Ollama \cite{ollama2024} running \texttt{nomic-embed-text} locally (768-dimensional vectors, no API key required).
  \item \textbf{Vector store}: ChromaDB \cite{chromadb2024} with persistent storage, cosine similarity, and metadata filtering by \texttt{user\_id}.
  \item \textbf{Collections}: Two collections---\texttt{user\_memories} (personality, interests, work, relationships) and \texttt{conversation\_history} (semantic search over past messages).
\end{itemize}

\subsubsection{Memory Lifecycle}

\begin{enumerate}[leftmargin=*]
  \item \textbf{Acquisition}: The LLM is instructed to call \texttt{remember\_about\_user} whenever the user shares personal information. The system prompt provides explicit examples: \textit{``if a user says `I'm a software developer and I like hiking,' save TWO memories: one for work and one for interests.''}
  \item \textbf{Embedding}: Each memory is embedded via Ollama and stored in ChromaDB with the user ID and category as metadata.
  \item \textbf{Retrieval}: On each new message, the user's query is embedded and the top-10 most relevant memories are retrieved by cosine similarity, filtered by user ID.
  \item \textbf{Injection}: Retrieved memories are appended to the system prompt with relevance scores, providing context without loading all memories.
\end{enumerate}

This approach scales to thousands of memories per user---only the relevant subset is retrieved per conversation turn, keeping the context window small.

\subsection{Vibe-Coding Workspace (P4)}

The workspace system provides a full git-based development environment for the agent to autonomously develop features.

\subsubsection{Tool Suite}

Fifteen workspace tools cover the complete development workflow:

\begin{itemize}[leftmargin=*]
  \item \textbf{Exploration}: \texttt{list\_files}, \texttt{read\_file}, \texttt{grep} (regex), \texttt{find} (glob), \texttt{inspect} (AST-based extraction of classes, functions, signatures, and docstrings)
  \item \textbf{Modification}: \texttt{write\_file}, \texttt{edit\_file} (precise find-and-replace), \texttt{delete\_file}, \texttt{install} (auto-detects pip/uv/npm)
  \item \textbf{Verification}: \texttt{run} (any shell command---pytest, npm run build, ruff), \texttt{check\_syntax} (py\_compile + ruff for Python, tsc for TypeScript)
  \item \textbf{Git}: \texttt{diff}, \texttt{commit\_push} (with auto-cleanup), \texttt{create} (clone + branch), \texttt{cleanup}
\end{itemize}

\subsubsection{Adaptive Workflow}

The system prompt defines a four-phase adaptive workflow:

\begin{enumerate}[leftmargin=*]
  \item \textbf{Understand}: Clone the repository, explore the file structure, read relevant code, grep for patterns, inspect function signatures.
  \item \textbf{Research}: Search online documentation for unfamiliar libraries using \texttt{web\_search} and \texttt{scrape\_fetch\_page}.
  \item \textbf{Implement}: Install dependencies, write/edit code, check syntax on each changed file.
  \item \textbf{Verify (loop)}: Run pytest and npm run build. If either fails, read the error output, fix the code, and re-run. Repeat until all checks pass. If stuck after three attempts, report the issue to the user.
\end{enumerate}

The workspace is automatically cleaned up after a successful \texttt{commit\_push}, and expires after one hour if abandoned.

\subsubsection{Code Inspection}

The \texttt{workspace\_inspect} tool uses Python's \texttt{ast} module to extract structured information from source files:

\begin{lstlisting}[language=Python, caption={AST-based code inspection output}]
{"classes": [{"name": "Agent", "bases": ["object"],
  "methods": [{"name": "__init__", "signature": "def __init__(self, provider: LLMProvider, ...)",
               "docstring": "Initialize the agent with a provider and system prompt."}]}],
 "functions": [...], "imports": [...], "constants": [...]}
\end{lstlisting}

This allows the agent to understand APIs and constructors before using them, rather than reading entire files.

\subsection{Self-Evolution: The Assistant That Extends Itself}

A distinctive feature of OpenPA is that the vibe-coding system can be directed at the assistant's \textit{own} codebase. Since OpenPA's source code is hosted on GitHub (\texttt{maxwellau2/OpenPA}), the same workspace tools that modify external repositories can modify OpenPA itself.

When a user says ``add a LinkedIn tool to OpenPA,'' the agent executes the following self-modification workflow:

\begin{enumerate}[leftmargin=*]
  \item \textbf{Clone self}: \texttt{ws\_workspace\_create(repo="maxwellau2/OpenPA", branch="feature/linkedin-tool")}
  \item \textbf{Learn patterns}: Read an existing tool file (e.g., \texttt{backend/tools/mastodon.py}) and the tool registry (\texttt{backend/tools/registry.py}) to understand the FastMCP tool pattern, credential handling, and namespace mounting convention.
  \item \textbf{Generate}: Create \texttt{backend/tools/linkedin.py} following the established pattern, with appropriate tool definitions, OAuth credential retrieval, and API calls.
  \item \textbf{Register}: Update \texttt{registry.py} to import and mount the new tool server.
  \item \textbf{Test}: Write unit tests in \texttt{backend/tests/test\_linkedin.py} and run the full test suite.
  \item \textbf{Build}: Run \texttt{npm run build} on the frontend to ensure no breakage.
  \item \textbf{Ship}: Commit, push, and create a pull request for human review.
\end{enumerate}

This makes OpenPA a \textbf{self-evolving platform}: the set of services it supports is not fixed at development time but can grow through natural language requests. The MCP-based tool architecture ensures that new tools are automatically discoverable by the agent loop once registered, without requiring changes to the agent, the system prompt, or the frontend.

\subsubsection{Evolution Triggers}

Self-evolution can be triggered through multiple channels:

\begin{itemize}[leftmargin=*]
  \item \textbf{Direct user request}: A user types ``add a LinkedIn tool'' in the chat interface. The agent enters workspace mode and executes the self-modification workflow immediately.
  \item \textbf{GitHub Issues}: The agent can monitor its own repository for new issues using \texttt{github\_list\_notifications}. When a user (or community contributor) files an issue such as ``Feature request: add Notion integration,'' the agent can read the issue, understand the requirements, and autonomously begin implementation---creating a branch, writing the code, and linking the resulting PR to the original issue.
  \item \textbf{Scheduled evolution}: Using the \texttt{scheduler\_schedule\_task} tool, periodic self-improvement cycles can be scheduled. For example, a nightly job could check for open feature-request issues and attempt to implement the highest-priority one.
\end{itemize}

\subsubsection{Safety Bounds}

The self-evolution capability is bounded by multiple safety mechanisms:

\begin{itemize}[leftmargin=*]
  \item Changes are submitted as \textbf{pull requests}, not merged directly---a human reviews before deployment.
  \item The \textbf{full test suite must pass} before pushing; broken code is never shipped.
  \item The workspace is \textbf{cleaned up automatically} after push, preventing resource leakage.
  \item The agent operates on a \textbf{cloned copy} of the repository, not the live deployment, eliminating the risk of corrupting the running instance.
\end{itemize}

\subsection{System Prompt Design}

The system prompt (approximately 4,500 words) follows a structured format inspired by Claude Code \cite{anthropic2025claudecode} and Pi \cite{badlogic2025pi}:

\begin{enumerate}[leftmargin=*]
  \item \textbf{Identity and scope}: Establishes the agent's role and capabilities.
  \item \textbf{Action rules}: Always use tools; never apologise for lack of access; parallelise independent tool calls.
  \item \textbf{Memory rules}: Proactively save user facts; check existing memories before acting.
  \item \textbf{Tool preference hierarchy}: Explicit mapping of task types to tool categories (e.g., ``use workspace tools for repo work, sandbox for one-off scripts'').
  \item \textbf{Output rules}: Be concise; lead with the answer; summarise tool output.
  \item \textbf{Coding rules}: Don't over-engineer; read before writing; never push broken code.
  \item \textbf{Workflow templates}: Detailed vibe-coding and self-evolution workflows.
  \item \textbf{Concrete examples}: 20+ example user queries mapped to expected tool call sequences.
\end{enumerate}

Dynamic context---including the user's email, connected services, and RAG-retrieved memories---is appended after the static rules, forming a ``prompt sandwich'' where stable instructions bracket variable context.

\subsection{Multi-Provider LLM Support (P5)}

OpenPA defines an abstract \texttt{LLMProvider} interface with a single method:

\begin{lstlisting}[language=Python]
async def chat(self, messages: list[Message],
               tools: list[dict] | None, system: str | None) -> LLMResponse
\end{lstlisting}

Six providers implement this interface, each converting the unified message and tool format to provider-specific schemas.

OpenPA adopts a \textbf{Bring-Your-Own-Key (BYOK)} model: the server stores no LLM provider credentials. Each user configures their own API keys through the settings page, and these are stored per-user in the database. This design has several advantages:

\begin{itemize}[leftmargin=*]
  \item \textbf{Cost transparency}: Users pay their own provider bills directly; the platform operator bears no inference cost.
  \item \textbf{No vendor lock-in}: One user may use Gemini's free tier for casual tasks while another uses Claude Sonnet for complex workflows---on the same OpenPA instance.
  \item \textbf{Privacy}: Users with strict data requirements can use Ollama \cite{ollama2024} with locally-hosted models (e.g., Gemma 4, Qwen 3.5, Llama 3.1), ensuring no data leaves their machine.
  \item \textbf{Scalability}: The server-side cost of adding users is near-zero since inference is offloaded to external providers or local hardware.
\end{itemize}
